% generated by R/50 (v2026a); csv: t1_convergent_validity.csv
\begin{table}[ht]
\caption{Convergent validity: correlation of the LLM scores with hand-coded CMP indices}
\label{tab:indep:cmpvalid}
\footnotesize
\centering
\resizebox{\ifdim\width>\linewidth\linewidth\else\width\fi}{!}{%
\begin{tabular}{lrrrrrrrrr}
\toprule
LLM score & CMP liberalism & CMP political & CMP social & CMP social (core) & CMP economic & CMP economic (market) & CMP economic (welfare) & CMP rile (reversed) & CMP populism \\
\midrule
Liberalism (overall) & 0.33 & $-$0.01 & 0.07 & 0.10 & 0.43 & 0.51 & 0.37 & $-$0.14 & $-$0.21 \\
Political liberalism & 0.33 & 0.13 & 0.20 & 0.19 & 0.20 & 0.29 & 0.15 & 0.07 & $-$0.21 \\
Social liberalism & 0.42 & 0.24 & 0.56 & 0.48 & $-$0.14 & $-$0.05 & $-$0.22 & 0.43 & $-$0.22 \\
Economic liberalism & 0.15 & $-$0.18 & $-$0.27 & $-$0.18 & 0.65 & 0.69 & 0.63 & $-$0.45 & $-$0.11 \\
Financial-market liberalism & 0.14 & $-$0.21 & $-$0.24 & $-$0.15 & 0.61 & 0.66 & 0.58 & $-$0.42 & $-$0.15 \\
Populism (overall) & $-$0.15 & 0.11 & 0.05 & 0.04 & $-$0.36 & $-$0.42 & $-$0.29 & 0.12 & 0.39 \\
\bottomrule
\end{tabular}}
\par\smallskip
\begin{minipage}{\linewidth}\scriptsize Pearson correlations at the manifesto level, pooled over all manifestos with both scores. Rows: LLM scores (cross-model means); columns: CMP indices, log(R + 0.5) - log(L + 0.5) of the quasi-sentence counts of the pro- and anti-liberal categories (Lowe et al. 2011); CMP rile reversed so that higher = more liberal.\end{minipage}
\end{table}
